% -------------------------------------------------------------------------------------
% Plantilla para escribir una tesis de la Universidad Nacional de Colombia en LaTeX
% *************************************************************************************
% -------------------------------------------------------------------------------------
\PassOptionsToPackage{spanish,english}{babel}
\documentclass[]{thesisUnal}
% -------------------------------------------------------------------------------------
% Espacio reservado para la carga de los paquetes por parte del autor
% **** --------------------------------------------------------------------------------

\usepackage[spanish]{babel}
\usepackage{rotating}
\usepackage{inputenx}
\usepackage[T1]{fontenc}
\usepackage{amsmath}
\usepackage{amssymb}
%\usepackage[spanish,activeacute,es-tabla]{babel}
%\usepackage[utf8]{inputenc}
\usepackage{amsmath, amsthm, amsfonts}
\usepackage{graphicx} % LaTeX
\usepackage{float} % LaTeX
\usepackage{bm}
%\usepackage{natbib} % bibliograf�a
%\usepackage{enumitem}


% //// --------------------------------------------------------------------------------
% -------------------------------------------------------------------------------------
% Espacio reservado para colocar las definiciones especiales por parte del autor
% **** --------------------------------------------------------------------------------




%%%%%%%%%%%%%%%%%%%%%%%%%%%%%%%%%%%%%%%%%%%%%%%%%%%%%%%%%%%%%%%%%%%%%%%%%%%%%%%%%%%%%%%
% \includeonly{cap1,cap2}
%%%%%%%%%%%%%%%%%%%%%%%%%%%%%%%%%%%%%%%%%%%%%%%%%%%%%%%%%%%%%%%%%%%%%%%%%%%%%%%%%%%%%%%
% //// --------------------------------------------------------------------------------
% -------------------------------------------------------------------------------------
% * CUERPO DEL DOCUMENTO
% **** --------------------------------------------------------------------------------
% -------------------------------------------------------------------------------------
\begin{document}

\logouniversity[45mm]{imagesThesis/logo_university_nacho}

\infothesis[
    author = {Sergio David Solano Bejarano},
    degreeauthor = {Ingeniero Industrial},
    %code = {000000},
    advisor = {B. Piedad Urdinola Contreras, Ph.D.},
    degreeadvisor = {Doctor en Demograf�a},
    %coadvisor = {Jonatan Gom�z Perdomo, Ph.D.},
    %degreecoadvisor = {Doctora en Demograf�a},
    title = {Proyecci�n de poblaciones carcelarias en Colombia},
    titledegree = {Disertaci�n presentada para optar al t�tulo de},
    degree = {Master en Ciencias - Estad�stica},
    researchline = {Demograf�a},
    %researchgroup = {\LaTeX: para el fomento del uso de \LaTeX\ en la investigaci�n},
    university = {Universidad Nacional de Colombia},
    faculty = {Facultad de Ciencias},
    department = {Departamento de Estad�stica},
    city = {Bogot�, D.C.},
    date = {Abril de 2017},
]

\abstractthesis[
    titlespanish = {Proyecci�n de poblaciones carcelarias en Colombia},
    abstractspanish = {Se realizaron proyecciones de la poblaci�n carcelaria en Colombia, usando la informaci�n disponible para los a�os 1991-2017. La informaci�n publicada periodicamente no incluye las tasas de transici�n (ingreso y salida) del sistema, por esta raz�n se eligieron tres m�todos que permiten realizar la proyecci�n a partir de la poblaci�n observada. Los m�todos utilizados son: modelos demogr�ficos para poblaciones peque�as, modelos ARIMA, y modelos Estado-Espacio. Se compar� el ajuste de cada modelo, sus ventajas y desventajas.},
	keywordspanish = {Poblaciones carcelarias, series de tiempo, procesos SARIMA, Modelos Estado Espacio, Poblaciones peque�as},
    titleenglish = {Prison populations projections for Colombia},
    abstractenglish = {Projections of the prison population in Colombia are made, considering available data from 1991-2017. Monthly released data does not include admission or release rates, so we use three methods that work over the total population. The considered methods are: Demographical methods for subnational populations, ARIMA Models, and State-Space Models. We compare the fit of the three models, their advantages and disadvantages.},
    keywordenglish = {Prison populations, time series, SARIMA processes, State Space Models, Subnational Populations},
]

\acceptationnote[
    note = {Aprobado},
    mention = {Meritoria o Laureada},
    jury = {Jurado uno},
    jury = {Jurado dos},
    %jury = {Michel Goossens},
    advisor = {B. Piedad Urdinola},
    %coadvisor = {Jonatan Gom�z},
    date = {Bogot�, D.C., Junio 01 de 2017}
]

\frontmatter
    \dedicatory{dedicatoria}
    \acknowledgement{agradecimientos}
    \tableofcontents
    \listoftables
    \listoffigures
    \introduction{introduccion}
\mainmatter
	\chapter{Antecedentes te�ricos}

En 2016 Colombia ocup� el puesto catorce entre doscientos cincuenta y un paises por el tama�o de su poblaci�n carcelaria (120 914 hbts.) y el cincuenta y uno seg�n la tasa de encarcelamiento (240 por cada 100.000 hbts). Tasa que pas� de 51,5 en  el a�o 2000 a 240 por cada 100.000 hbts en 2016. Con una ocupaci�n del 154\% de las plazas disponibles, resulta relevante contar con proyecciones de la poblaci�n carcelaria en el corto, mediano y largo plazo.\cite{InstitueforCriminalPolicyResearch2016}

\section{Proyecciones de poblaci�n}

``Una estimaci�n poblacional consiste en determinar el tama�o o las caracter�sticas de una poblaci�n, para el momento actual o para uno anterior, en ausencia de informaci�n. Cuando se realizan un conjunto de supuestos sobre el comportamiento de los vitales hacia el futuro, hablamos de proyecci�n, y cuando se escoge un escenario como el m�s probable, hablamos de pron�stico'' \cite {Swanson2012}.

Para incluir la incertidumbre en las proyecciones de poblaci�n Lee(1994) considera los siguientes m�todos \cite {Lee1994}: 

\begin{itemize}
	\item El enfoque de escenarios alto, medio y bajo: Asume comportamientos fijos para la fertilidad, la mortalidad y las migraciones durante el per�odo de proyecci�n, basado en algunos supuestos \cite {Lee1994}.
	\item An�lisis estoc�sticos
	\begin{itemize}
		\item An�lisis ex-post: consiste en evaluar el error de pron�stico en proyecciones anteriores y aplicarlo a las nuevas proyecciones \cite {Lee1994}.
		\item Simulaci�n estoc�stica: Permite hacer proyecciones de poblaci�n, al asignar una distribuci�n de probabilidad a las tasas vitales (mortalidad, natalidad, migraciones) \cite {Lee1994}.
		\item Modelos estoc�sticos de la tasa de crecimiento: Consiste en estimar la tasa de crecimiento del total de la poblaci�n; aunque permite estimar intervalos de confianza, no permite separar la proyecci�n de las tasas vitales, ni de las franjas etarias \cite {Lee1994}.
		\item Matrices de Leslie con modelos estimados para las tasas vitales:  Al usar matrices de Leslie se estima la poblaci�n por rangos etarios para un instante i, y se calcula la poblaci�n en el instante i + 1 aplicando la natalidad y la mortalidad proyectadas para el per�odo i. Puesto que las series de poblaci�n carcelaria no se publican separadas por edad, no se abordar� esta t�cnica.
	\end{itemize}
\end{itemize}

\subsection{Proyecci�n de poblaciones peque�as}

La proyecci�n de �reas peque�as es entendida como la proyecci�n a un nivel geogr�fico menor al nacional. Estas proyecciones pueden incluir, departamentos, ciudades o poblaciones especiales \cite {Swanson2012}.
``Una poblaci�n especial es un grupo poblacional que se encuentra restringido a un �rea por una medida administrativa o legislativa. Dentro de los grupos usualmente considerados se encuentran las prisiones, universidades, hospitales e instituciones militares''.
Este tipo de poblaci�n puede tener una estructura etaria y de sexo, y unos vitales diferentes al resto de la poblaci�n; adem�s no suelen envejecer en el mismo lugar, lo que permite mantener una estructura etaria que no var�a a trav�s del tiempo. \cite {Swanson2012}.

\subsection{Aplicaciones nacionales e internacionales}

Las proyecciones de poblaciones carcelarias oficiales analizadas corresponden, en buena parte, a los m�todos expuestos en los cap�tulos anteriores:  Proyecciones por escenarios, proyecci�n de la tasa de crecimiento, modelos ARIMA para las tasas de ingreso y salida.

En Colombia (CONPES 3828) se proyect� la poblaci�n carcelaria usando la tasa media de crecimiento anual (1993-2014) \cite {DepartamentoNacionaldePlaneacion2015}. Esta proyecci�n no tiene en cuenta la incertidumbre asociada con las variaciones aleatorias en las tasas, ni las asociadas a cambios estructurales, ni que son una poblaci�n especial, con patrones muy distintos a los del resto de la poblaci�n \cite{Swanson2012} .

El Reino Unido hasta 2015 realizaba una proyecci�n por escenarios (alto, medio y bajo), a�o en el cual cambi� a un modelo de proyecci�n de la media y su incertidumbre. La incertidumbre se incluy� a trav�s de un an�lisis ex-post, de la  desviaci�n de la proyecci�n en a�os anteriores \cite {Justice2014}.

El departamento de Justicia de los Estados Unidos realiz� estimaciones de la poblaci�n carcelaria por estado para el per�odo 2013-2014. Las estimaciones parten del censo de prisiones 1993-2014. Estas proyecciones se puede enmarcar dentro de las proyecciones de �reas peque�as \cite {Minton2015}.

El bureau de estad�sticas e investigaci�n del crimen en Australia proyecta las tasas de arresto y sentencia usando modelos ARIMA;  a partir de estas tasas proyecta la poblaci�n carcelaria. Estas proyecciones incluyen un per�odo de validaci�n de tres a�os. Los resultados mostraban que la serie real se encuentra dentro de los intervalos de confianza de la proyecci�n, cercano a la proyecci�n de la media \cite{Wan1}.

Blummstein desarrolla un m�todo de proyecci�n basado en los componentes demogr�ficos, tasas espec�ficas de arresto por delito y reincidencias, a partir de estos datos proyecta el tama�o y la composici�n de las poblaciones \cite{Blumstein1980a}.

Con los datos libres disponibles en Colombia no se podr�a utilizar el enfoque ARIMA ni el m�todo de Blummstein, pues las tasas de encarcelamiento y sentencia no se publican. La tesis busca proponer un m�todo de proyecci�n, para situaciones donde no se cuenta con el registro de los vitales o su equivalente en la poblaci�n analizada, en este escenario resulta conveniente usar las series ARIMA sobre la series de poblaci�n sindicada y poblaci�n sindicada, y los modelos estado-espacio para la serie bivariada. Se proponen los m�todos demogr�ficos de estimaci�n para �reas peque�as Censal-ratio y Ratio-Correlation, y su contraste con las metodolog�as ARIMA y Estado-Espacio que a diferencia de los m�todos demogr�ficos mencionados incorporan una medida de incertidumbre.

\section{Series de tiempo}

Una serie de tiempo es un conjunto de observaciones ${y_t}$ asociadas a un instante de tiempo t. Es usual referirse como series de tiempo, tanto a las realizaciones ${y_t}$ como a las variables aleatorias ${Y_t}$ que las generan.\cite{Brockwell2011} 

En una regresi�n lineal cl�sica, una variable $Y$ es explicada o predicha en funci�n de un conjunto de n variables $X$. La diferencia entre el valor observado en el instante t $y_t$ y el valor predicho se suponen provenientes de un proceso aleatorio con media cero. \cite{Commandeur2007}

\begin{equation}\label{eq:regresion}
y_t = \beta_{0} + \beta_{1} x_{t1} + \beta_{2} x_{t2} ... \beta_{n} x_{tn}  + w_t
\end{equation}

Tambi�n es posible representar la observaci�n $y_t$ en funci�n de las observaciones anteriores como:

\begin{equation}\label{eq:serierecursiva}
y_t = \beta_{0} + \beta_{1} y_1 + \beta_{2} y_2 + ... + \beta_{t-1} y_{t-1} +  w_t
\end{equation}

La variable aleatoria $Y$ en el instante i, es una combinaci�n lineal de los valores observados de la variable aleatoria, m�s un error aleatorio, donde $\{\epsilon_1, \epsilon_2, ...\}$ son independientes e id�nticamente distribuidas (i.i.d.). \cite{Shumway} Este supuesto no suele cumplirse, raz�n por la cual el an�lisis de series de tiempo se ha desarrollado como un �rea particular de la estad�stica \cite{Brockwell2011}.

Se define un proceso como ruido blanco a la colecci�n de variables aleatorias $z_t$, no correlacionadas, con media $0$ y varianza $\sigma^2_w$. \cite{Shumway}

Bajo un modelo de $k$ par�metros, el estimador m�ximo veros�mil de la varianza es: 

\begin{equation}\label{varianza}
\hat {\sigma}^2_k = \frac{SSE(k)}{n} = \frac{1}{n} \sum_{t = 1}^{n}  (y_t - \hat{y_t})
\end{equation}

Para seleccionar el modelo que mejor explique los datos observados se recurre a los criterios AIC (Criterio de Informaci�n de Aikake) y BIC (Criterio de Informaci�n Bayesiano), que penalizan los modelos con mayor cantidad de par�metros y est�n definidos como: 

\begin{equation}\label{AIC}
AIC = ln (\hat{\sigma_k}^2) + \frac{n+2k}{n}
\end{equation}

\begin{equation}\label{AIC}
BIC = ln (\hat{\sigma_k}^2) + \frac{k (ln(n))}{n}
\end{equation}

donde $k$ es la cantidad de par�metros en el modelo y $n$ es la cantidad de observaciones \cite{Tsay2002}.  Tanto en el AIC como en el BIC se selecciona el modelo de $k$ par�metros con menor valor.  

Para medir la dependencia entre las observaciones recurrimos a la funci�n de autocovarianca y a la funci�n de autocorrelaci�n. La funci�n de autocovarianza se define como \cite{Shumway}: 

\begin{equation}\label{covarianza}
\gamma(s,t) = cov(y_{s}, y_{t}) = E[(y_{s} - \mu_{s})(y_{t} - \mu_{t})]
\end{equation}

para todo s y t, donde $\mu = EY_t$ y k representa la cantidad de pasos entre dos observaciones. La autocovarianza mide la dependencia lineal entre dos puntos en la misma serie en diferentes instantes. \cite{Shumway}. 

La funci�n de autocorrelaci�n (ACF) se define como \cite{Brockwell2011}: 

\begin{equation}\label{correlacion}
\rho (s,t) = \dfrac{\gamma(s,t)}{\sqrt{\gamma(s,s) \gamma(t,t)}} 
\end{equation}


Para poder realizar predicciones sobre el estado futuro de una serie de tiempo, es necesario suponer que el comportamiento de la serie es estable a trav�s del tiempo, incluso con un componente aleatorio. La estacionaridad es el concepto que permite articular esta necesidad. Estamos ante un \textbf{proceso estacionario} cuando las variables $\{X_1,..., X_k\}$ tienen la misma distribuci�n conjunta que $\{X_{h+1},...,X_{h+k}\}$, para todos los enteros $h$ y $k$ \cite{Brockwell2011}.

La estacionaridad d�bil se presenta cuando $E[Y_t]$ y $E[Y_t,Y_{t+h}]$ son independientes de $t$, es decir: i) Presenta media $\mu$ constante e independiente de $t$ ii) La funci�n de autocovarianza $\gamma(h, h+k)$ depende solamente de la cantidad de pasos que separa las observaciones ($k$). En adelante, al referirse a estacionaridad, se tratar� de estacionaridad debil, a menos que se indique lo contrario.\cite{Shumway}

Cuando la serie no tiene un comportamiento estacionario se puede recurrir a diferenciar la serie, este proceso al ser de uso frecuente tiene su propia notaci�n. La diferencia $\bigtriangledown y_t = y_t - y_{t-1}$ se conoce como diferencia de primer orden. La diferencia de segundo orden se define como $\bigtriangledown^2 y_t = (y_t - y_{t-1}) - (y_{t-1} - y_{t-2} )$.

El operador de rezago B se define como \cite{Shumway}:

\begin{equation}\label{backshift}
By_t = y_{t_1} 
\end{equation}

tal que: 
\begin{equation}\label{backshift2}
\bigtriangledown y_t = (1-B) y_t 
\end{equation}

La notaci�n se puede extender de forma que: 

\begin{equation}\label{backshift3}
\bigtriangledown^d y_t = (1-B)^d y_t 
\end{equation}

La funci�n de auto-correlaci�n (ACF) de una serie de tiempo estacionaria  se define como \cite{Shumway}: 

\begin{equation}\label{autocorrelacion}
\rho (h) = \frac{\gamma (h)}{\gamma(0)}
\end{equation}

donde $\gamma (h) = cov(y_{t+h},y_t) = E[(y_{t-h} - \mu)(y_t- \mu)]$ y $\bar{y}$ denota la media del proceso.


Para estimar la ACF utilizamos la funci�n de autocorrelaci�n muestral definida como \cite{Shumway}: 

\begin{equation}\label{autocorrelacionmuestral}
\hat{\rho (h)} = \frac{\hat{\gamma} (h)}{\hat{\gamma} (0)} = \frac{\sum_{t = 1}^{n - h}(y_{t+h}-\bar{y})(y_t - \bar{y})}{\sum_{t = 1}^{n}(y_t-\bar{y})^2}
\end{equation}
con $ h = 0,1, ..., n-1$.

Si tenemos las variables aleatorias $X,Y,Z$ y queremos medir la correlaci�n entre las variables $Z$ y $X$, descontando el efecto de la variable Y, recurrimos a la funci�n de correlaci�n parcial $\rho_{XY|Z}$, que se calcula como \cite{Shumway}: 

\begin{equation}\label{correlacionparcial}
\rho_{XY|Z} = corr\{ X-\hat{X},Y-\hat{Y}\}
\end{equation}

Al adaptar este concepto a series de tiempo tenemos que la funci�n de autocorrelaci�n parcial de un proceso estacionario $y_t$, $\alpha_{hh}$, con $h = 1, 2, ... , t$ se define como \cite{Shumway}: 

\begin{equation}\label{correlacionparcialestacionario1}
\alpha_{11} = corr(y_1,y_0)  = \rho(1)
\end{equation}

y 

\begin{equation}\label{correlacionparcialestacionario2}
\alpha_{hh} = corr(y_h - \hat{y_h},y_0 - \hat{y_0}), h \geq 2 
\end{equation}

donde $\hat{y_h}$ es la regresi�n de $y_h$ en $\{y_1,y_2, ..., y_{h-1}\}$  y $\hat{y_0}$ es la regresi�n de $y_0$ en $\{y_1,y_2, ..., y_{h-1}\}$.

En un proceso estacionario $x_t $ la PACF es la correlaci�n entre $ x_t,x_{t+h} $ una vez controlado por los elementos entre ellos $ \{x_{t-h+1}, x_{t-h+2},..., x_{t-1} \} $ \cite{Shumway}.

\subsection{Modelos SARIMA}

Box \& Jenkins proponen una aproximaci�n iterativa de  cuatro etapas para la selecci�n de un modelo \cite{Box2013}: 

\begin{enumerate}
	\item Selecci�n de una clase de modelos, con base en la teor�a y la pr�ctica.
	\item Identificaci�n del modelo, donde se seleccionan un conjunto de par�metros que permitan explicar el sistema con parsimonia. 
	\item Estimaci�n de par�metros. 
	\item Chequeo diagn�stico, para detectar fallas en el ajuste. Si se detectan fallas en el ajuste, se regresa al segundo paso. 
\end{enumerate}

Shumway en \cite{Shumway} plantea que un proceso es auto-regresivo de orden p o AR(p) si:

\begin{equation}\label{ar}
y_t = \phi_1 y_{t-1} + \phi_2 y_{t-2} ... + \phi_p y_{t-p} + w_t
\end{equation}

Donde $y_t$ es estacionaria con $\mu = 0$ ;  $ \phi_1,  \phi_2,...  ,\phi_p$ son constantes  ($\phi_p \neq 0$) y los $w_t$ son independientes e id�nticamente distribuidos. Cuando la $\mu \neq 0$

\begin{equation}\label{ar1}
y_t - \mu  = \phi_1 (y_{t-1} - \mu) + \phi_2 (y_{t-2}-\mu) ... + \phi_p (y_{t-p}-\mu) + w_t
\end{equation}

Que se puede representar como:

\begin{equation}\label{ar2}
y_t  = \mu(1 - \phi_1 ... - \phi_p) + \phi_1 (y_{t-1}) + \phi_2 (y_{t-2}) ... + \phi_p (y_{t-p}) + w_t
\end{equation}

Tambi�n es posible utilizar el operador B para representar la serie \ref{ar} como:

\begin{equation}\label{ar3}
y_t  - \phi_1 (y_{t-1}) - \phi_2 (y_{t-2}) ... - \phi_p (y_{t-p}) =  w_t
\end{equation} 

\begin{equation}\label{ar4}
(1 - \phi_1 B - \phi_2 B^2 ... - \phi_p B^p)y_t =  w_t
\end{equation} 

o en una notaci�n simplificada: 

\begin{equation}\label{ar5}
\phi(B)y_t =  w_t
\end{equation} 

Se considera un proceso de promedio m�vil de orden q o MA(p) si:

\begin{equation}\label{ma}
y_t = \theta_1 w_{t-1} + \theta_2 w_{t-2} ... \theta_q w_{t-q} + w_t
\end{equation}

Donde $ \theta_1,  \theta_2,...  ,\theta_q$ son par�metros  ($\phi_p \neq 0$) y los $w_t$ son independientes e id�nticamente distribuidos y $y_t$ tiene media cero. A diferencia de los procesos auto-regresivos, los procesos de promedio movil son estacionarios sin importar el valor de $\theta_1,  \theta_2,...  ,\theta_q$ \cite{Shumway}. Este proceso tambi�n se puede representar usando el operador de rezago al ponerlo en la forma: 

\begin{equation}\label{ma1}
y_t = (\theta_1 B + \theta_2 B^2 ... \theta_q B^q + 1) w_t
\end{equation}

o en forma simplificada: 

\begin{equation}\label{ma2}
y_t = \theta (B) w_t
\end{equation}

Un proceso se considera ARMA(p,q) si $y_t$ es estacionaria y  \cite{Shumway} : 

\begin{equation}\label{arma1}
y_t = \phi_1 y_{t-1} + \phi_2 y_{t-2} + ... + \phi_p y_{t-p} + \theta_1 w_{t-1} + \theta_2 w_{t-2} ... \theta_q w_{t-q} + w_t
\end{equation}

Con $E(y_t) = \mu = 0$. Cuando $(\mu \neq 0)$ se puede representar como: 

\begin{equation}\label{arma2}
y_t  = \mu(1 - \phi_1 ... - \phi_p) + \phi_1 y_{t-1} + \phi_2 y_{t-2} + ... + \phi_p y_{t-p} + \theta_1 w_{t-1} + \theta_2 w_{t-2} ... \theta_q w_{t-q} + w_t
\end{equation}

A este modelo se le conoce como ARMA(p,q). Usando las ecuaciones \ref{ar5} y  \ref{ma2} se puede representar como: 

\begin{equation}\label{arma3}
\phi(B)y_t  =  \theta (B) w_t
\end{equation}

Una t�cnica com�n para tratar con series no estacionarias consiste en diferenciarlas hasta que la serie obtenida sea ARMA(p,q) \cite{Brockwell2011}. Si la serie se diferencia $d$ veces el proceso se nota como ARIMA(p,d,q).

Un proceso ARIMA (i,1,j) resulta al considerar una serie de la forma \cite{Shumway}: 

\begin{equation}\label{ARIMA}
y_t = \alpha +  y_{t-1}
\end{equation}

Tal que el proceso $y_t - y_{t-1}$ es un proceso ARMA(p,q).

Los criterios para identificar el orden de un ARMA son \cite{Boland2011}: 

\begin{itemize}
	\item Cuando la funci�n de autocorrelaci�n (ACF) se reduce progresivamente, y  la funci�n de autocorrelaci�n parcial (PACF) no tiene picos en lags luego de p, es un proceso autoregresivo de orden p  AR(p) 
	\item Cuando la funci�n de autocorrelaci�n tiene un pico en el lag q, y la funci�n de autocorrelaci�n parcial se reduce progresivamente, es un proceso de media mmovil de orden q MA(q)
	\item Si ambas funciones se reducen gradualmente, se trata de un ARMA (p,q)
\end{itemize}

Cuando las series presentan un comportamiento estacional se recurre al proceso de diferenciaci�n estacional de la forma: \cite{Tsay2002}

\begin{equation}\label{ARIMA}
\bigtriangledown_s y_t = y_t - y_{t-s} = (1 - B^s) y_t
\end{equation}

Un proceso autoregresivo de orden $P$, promedio m�vil de orden $Q$ estacional $ARMA(P,Q)_s$, donde s corresponde a la periodicidad de la serie se define como \cite{Shumway}: 

\begin{equation}\label{sarma1}
\Phi_P  (B^s) y_t  =  \Theta_Q (B^s) w_t
\end{equation}

donde $\Phi_P (B^s) = 1- \Phi_1 B^s - \Phi_2 B^{2s}- ... - \Phi_P B^{Ps}$ se denomina operador estacional autoregresivo de orden P y $\Theta_Q (B^s) = 1 + \Theta_1 B^s + \Theta_2 B^{2s} + ... + \Theta_Q B^{Qs}$ se denomina operador estacional promedio movil de orden Q \cite{Shumway}.

Un proceso SARIMA(p,d,q)(P,D,Q) se define como: 

\begin{equation}\label{sarma1}
\Phi_P  (B^s) \phi_p (B^s) \bigtriangledown^D \bigtriangledown^d y_t  =  \Theta_Q (B^s) \theta_q (B) w_t
\end{equation}

donde $w_t$ es un proceso gausiano de ruido blanco.

El paquete $astsa$ se ha utilizado para estimar los modelos ARIMA en el capitulo 3. La funci�n $sarima$ del paquete $astsa$ es un manejador de la funci�n $arima$ del paquete base. \cite{Stoffer2014}. La funci�n $arima$ utiliza la funci�n $optim$ para encontrar estimadores de m�xima verosimilitud.\cite{RCoreTeam2017}.

El paquete $astsa$ se ha utilizado para realizar los pron�sticos en el capitulo 3. La funci�n $sarima.for$ del paquete $astsa$ es un manejador de la funci�n $predict.Arima$ del paquete base. \cite{Stoffer2014}. La funci�n $predict.Arima$ utiliza la funci�n $KalmanForecast$. En la siguiente secci�n se hablar� en detalle del filtro de Kalman. \cite{RCoreTeam2017}

Una vez se ha realizado la estimaci�n del modelo, se analiza el comportamiento del error. Para descartar autocorrelaci�n en el error se utiliza la estad�stica Q de Ljung, definida como:

\begin{equation}\label{ljung}
Q = n(n+2) \sum_{n=1}^{H} \frac{ \hat {\rho_{e}}^2(h)}{n-1}
\end{equation}

donde $H$ es el rezago analizado, $\hat{\rho_{e}}(h)$  es la autocorrelaci�n estimada del error en el rezago h. La hip�tesis nula es que los errores son no correlacionados. Bajo la hip�tesis nula $Q \sim \chi_{H}^2$. \cite{Ljung1978} \cite{Shumway}

\subsection{Modelos Estado-Espacio}

Los modelos Estado-Espacio en su forma m�s simple, emplea un vector autoregresivo de la forma \cite{Shumway2006}: 

\begin{equation}\label{StateSpace1}
\boldsymbol{y_t} = \Phi \boldsymbol{y_{t-1}} + \boldsymbol{w_t}
\end{equation}

donde $\boldsymbol{y_t}$ es un vector de $p$ componentes. $ \boldsymbol w_t$ es un vector de $p \times 1$ variables aleatorias independientes e id�nticamente distribuidas normalmente con media cero y matriz de covarianzas $ \boldsymbol Q$. Asumimos que el proceso es normal con un vector $ \boldsymbol {y_0}$ con media $ \boldsymbol {\mu_0}$ y matriz de covarianzas $ \boldsymbol \sum$. \cite{Shumway2006}


\begin{equation}\label{StateSpace2}
\boldsymbol{z_t} = A_t \boldsymbol{y_{t}} + \boldsymbol{v_t}
\end{equation}

donde $A_t$ es de tama�o $q \times p$. Se suponen $\boldsymbol {w_t}$ y $\boldsymbol {v_t}$ no correlacionados.

Variables ex�genas o inputs fijos pueden ingresar en los estados o las observaciones. En este caso tendremos un vector $u_t$ de tama�o $r \times 1$: 

\begin{equation}\label{StateSpace6}
\boldsymbol{y_t} = \Phi \boldsymbol{y_{t-1}} + \boldsymbol{\Upsilon} \boldsymbol{u_t} + \boldsymbol{w_t}
\end{equation}


\begin{equation}\label{StateSpace7}
\boldsymbol{z_t} = A_t \boldsymbol{y_{t}} + + \boldsymbol{\Gamma} \boldsymbol{u_t} + \boldsymbol{v_t}
\end{equation}

con $\boldsymbol{\Upsilon}$ y $\boldsymbol{\Gamma}$ de dimensi�n $p x r$ y $q x r$, respectivamente. 

El objetivo es tener estimaciones de $x_t$ a partir de las observaciones $y_1, y_2, ..., y_s$. Cuando $s < t$ se trata de predicci�n. Cuando $s = t$ se habla de filtrado y para $s > t$ de suavizado \cite{Shumway2006}.

\begin{equation}\label{StateSpace3}
\boldsymbol{y_t^s} = E\{\boldsymbol{y_t} |\boldsymbol{ z_s}\}
\end{equation}

\begin{equation}\label{StateSpace4}
\boldsymbol{P_{t}^s} = E\{\boldsymbol{(y_{t1}-y_{t1}^s)} |\boldsymbol{{(y_{t2}-y_{t2}^s)}}'\}
\end{equation}


    \chapter{La poblaci�n carcelaria en Colombia 1991 - 2017}

\section{Descripci�n de los datos}

\section{El sistema carcelario en Colombia}

\section{Implicaciones en el modelado}
    \chapter{Modelos SARIMA}

La estimaci�n de los modelos se realiza usando el paquete base del software R \cite{RCoreTeam2017}, y el paquete astsa \cite{Stoffer2014}, cuyo uso es discutido en detalle por Shumway \cite{Shumway}. 

En adelante nos referiremos a la funci�n de auctocorrelaci�n como ACF y a la funci�n de autocorrelaci�n parcial como PACF. 

\section{Identificaci�n del modelo}

% Variaci�n poblaci�n total, sindicada y condenada
\begin{figure}[H]
	\centering
	\includegraphics[width=0.7\linewidth]{variacion_intermensual}
	\caption[Variaci�n inter-mensual de poblaci�n carcelaria]{Variaci�n inter-mensual de poblaci�n carcelaria, sindicados y condenados}
	\label{fig:variacion_intermensual}
\end{figure}

Antes de realizar la proyecci�n de una serie de tiempo, es necesario identificar el modelo que explique adecuadamente su comportamiento. 

Aunque conocemos, por la estructura del proceso que genera los datos, que las series de tiempo de poblaci�n sindicada y condenada no son independientes, podemos simplificar la proyecci�n, trat�ndolas como independientes. En este caso, los modelos ARIMA y SARIMA resultan apropiados, pues permiten explicar separadamente cada observaci�n en funci�n del comportamiento hist�rico de la serie.

El cap�tulo anterior suger�a que la poblaci�n carcelaria, tanto sindicada como condenada, tiene una marcada tendencia al alza. En este caso una herramienta �til es mostrar gr�ficamente la variaci�n mes a mes de la poblaci�n. \ref{fig:variacion_intermensual}. Para tener una mejor estrucutura de an�lisis usamos la funci�n "decompose" de R base. 

La serie tiene un componente estacional marcado, con una reducci�n de la poblaci�n carcelaria en diciembre. La variabilidad del componente aleatorio es elevada. La tendencia parece tener cambios estructurales en algunos periodos, por ejemplo reducci�n de la poblaci�n carcelaria entre 2005-2007,  y 2012 - 2015, e incrementos de la poblaci�n de magnitud mayor al promedio entre 2008 y 2012.

La mayor parte de la variaci�n de la poblaci�n total se puede asociar con variaciones en a poblaci�n sindicada \ref{fig:variacion_mensual_sindicados_desc}. La poblaci�n condenada tiene una tendencia con una tendencia m�s estable.  \ref{fig:variacion_mensual_condenados_desc}

% Descomposici�n variaci�n poblaci�n total
\begin{figure}[H]
\centering
\includegraphics[width=0.7\linewidth]{variacion_mensual_total_desc}
\caption[Descomposici�n de la variaci�n inter-mensual de poblaci�n carcelaria total]{Variaci�n mensual de poblaci�n carcelaria descompuesta por tendencia, estacionalidad y componente aleatorio.}
\label{fig:variacion_mensual_total_desc}
\end{figure}

% Descomposici�n variaci�n poblaci�n sindicada
\begin{figure}[H]
	\centering
	\includegraphics[width=0.7\linewidth]{variacion_mensual_sindicados_desc}
	\caption[Descomposici�n de la variaci�n inter-mensual de poblaci�n carcelaria total]{Variaci�n mensual de poblaci�n carcelaria sindicada,  descompuesta por tendencia, estacionalidad y componente aleatorio.}
	\label{fig:variacion_mensual_sindicados_desc}
\end{figure}

% Descomposici�n variaci�n poblaci�n condenada
\begin{figure}[H]
	\centering
	\includegraphics[width=0.7\linewidth]{variacion_mensual_condenados_desc}
	\caption[Descomposici�n de la variaci�n inter-mensual de poblaci�n carcelaria total]{Variaci�n mensual de poblaci�n carcelaria condenada,  descompuesta por tendencia, estacionalidad y componente aleatorio.}
	\label{fig:variacion_mensual_condenados_desc}
\end{figure}


Se trabaja sobre la diferencia de cada serie de tiempo con lag 1 y lag 12, al observar que la serie tiene un crecimiento sostenido entre 1991 y 2017, sugiriendo un proceso estocastico integrado. Sobre cada serie (poblaci�n total, poblaci�n sindicadda y poblaci�n condenada)  se realiza la funci�n de autocorrelaci�n y autocorrelaci�n parcial y se presenta en la figura \ref{fig:ACF_variacion_mensual}. Estas funciones son usadas como herramienta de diagn�stico, para intuir modelos adecuados en cada serie.

Con base en la tabla \ref{ARMA_ACF} se realiza una revisi�n del comportamiento de las series de poblaci�n carcelaria, poblaci�n sindicada y poblaci�n condenada.

% Pendiente traducir este fragmento
\begin{table}[H]
	\centering
	\caption{Comportamiento de la ACF y PACF en modelos ARMA(p,q) \cite{Shumway}}
	\label{ARMA_ACF}
	\begin{tabular}{lllll}
		\multicolumn{1}{|l|}{}              & \multicolumn{1}{c|}{\textbf{AR(p)}}       & \multicolumn{1}{c|}{\textbf{MA(q)}}       & \multicolumn{1}{c|}{\textbf{ARMA(p,q)}} &  \\ \cline{1-4}
		\multicolumn{1}{|l|}{\textbf{ACF}}  & \multicolumn{1}{l|}{Tails off}            & \multicolumn{1}{l|}{Cuts off after lag q} & \multicolumn{1}{l|}{Tails off}          &  \\ \cline{1-4}
		\multicolumn{1}{|l|}{\textbf{PACF}} & \multicolumn{1}{l|}{Cuts off after lag p} & \multicolumn{1}{l|}{Tails off}            & \multicolumn{1}{l|}{Tails off}          &  \\ \cline{1-4}
	\end{tabular}
\end{table}


\begin{figure}[H]
	\centering
	\includegraphics[width=0.7\linewidth]{ACF_var_pob}
	\caption[Autocorrelaci�n parcial de la variaci�n inter-mensual]{Autocorrelaci�n parcial de la variaci�n inter-mensual de la poblaci�n}
	\label{fig:ACF_variacion_mensual}
\end{figure}

\subsection{Poblaci�n total}

La tabla \ref{ARMA_ACF} presenta una gu�a para la interpretaci�n de la ACF y la PACF. En la \ref{fig:ACF_variacion_mensual} podemos observar que tanto la ACF como la PACF decaen lentamente. La APCF decae progresivamente en los meses 12, 24, 36, mientra que la ACF se corta en el lag 24, lo que sugiere que el proceso es un promedio m�vil de orden 1 en el componente estacional.

En este caso una primera estimaci�n se realiza para la proyecci�n de la poblaci�n total como SARIMA(1,1,1,0,0,1)


\subsection{Poblaci�n sindicada}

En la \ref{fig:ACF_variacion_mensual} podemos observar que la ACF decae lentamente, mientras la PACF cae por debajo del error luego del lag 1. La APCF decae progresivamente en los meses 12, 24, 36, mientra que la ACF se corta en el lag 12, lo que sugiere que el proceso es un promedio m�vil de orden 1 en el componente estacional.

En este caso una primera estimaci�n se realiza para la proyecci�n de la poblaci�n sindicada como SARIMA(1,1,1,0,0,1)

\subsection{Poblaci�n condenada}

El comportamiento de la ACF y la PACF es similar al de la poblaci�n sindicada, sugiriendo el mismo modelo: SARIMA(1,1,1,0,0,1).

\section{Estimaci�n de par�mtetros}
El proceso de estimaci�n se realizar� para cada serie por separado, luego se validar� la efectividad de incluir m�s par�metros, usando como criterio de comparaci�n el BIC. A modo informativo se usa la funci�n auto.arima del paquete forecast, para seleccionar el modelo con menor AIC. \cite{Hyndman2017}  
 \cite{Hyndman2008} 
 
\subsection{Poblaci�n total}

\subsection{Poblaci�n sindicada}

\subsection{Poblaci�n condenada}


\section{Proyecci�n}

\subsection{Poblaci�n total}

\subsection{Poblaci�n sindicada}

\subsection{Poblaci�n condenada}

\section{Conclusiones}

Al usar un ARIMA
    \chapter{M�todos demogr�ficos para poblaciones peque�as}

No todas las personas tenemos la misma probabilidad de ser encarceladas. Los hombres tienen tasas de encarcelamiento m�s altas que las mujeres, y las personas entre 18 y 40 a�os son m�s susceptibles que los mayores de 40 \cite{Blumstein1980a}. Las cifras  de encarcelamiento, presentadas en la figura \ref{fig:inpec_edad_genero1} sugieren que estas tasas son estables en periodos cortos de tiempo y proporcionales entre los diferentes rangos de edad (si el encarcelamiento en un rango de edad se eleva sucede un incremento proporcional en los otros).  

La poblaci�n colombiana ha cambiado significativamente desde 1991, tanto en su magnitud como en su composici�n por edad. Hemos pasado de 35 millones en 1991 a 45 millones en la segunda d�cada del siglo XXI. Tambi�n hemos presenciado una reducci�n en la proporci�n de menores de 18 a�os, que ha pasado de XXX en 1991 a XXX en 2016. Se proyecta que en 2020 tendremos tantas personas entre 15 y 29 a�os, como entre 0 y 14, cuando en 1991 la diferencia era del XXX \%. \cite{DANEPROY}

El grupo etario con mayor tasa de encarcelamiento est� entre los 18 y 29 a�os.  \cite{Blumstein1980a}. La poblaci�n total en este grupo crece sostenidamente en el periodo 1991-2020. Como porcentaje de la poblaci�n total, decrecen en dos puntos porcentuales en el mismo periodo, cediendo espacio a una poblaci�n de mayor edad. El crecimiento de este grupo etario, particularmente en hombres podr�a ser un buen predictor de la poblaci�n carcelaria.

Para incorporar la relaci�n entre la poblaci�n nacional y la poblaci�n carcelaria se recurre al m�todo censal-ratio,com�nmente usado en la proyecci�n de poblaciones peque�as. Este m�todo requiere la inclusi�n de variables sintom�ticas que agreguen informaci�n sobre la poblaci�n peque�a.

\begin{figure}[H]
	\centering
	\includegraphics[width=0.9\linewidth]{POB_NAL_EDAD-1.pdf}
	\caption[Proyecci�n de la poblaci�n nacional por rango de edad]{Proyecci�n de la poblaci�n nacional por rango de edad}{Fuente: DANE} {Elaboraci�n propia}
	\label{fig:pob_nal_dane_edad}
\end{figure}


\section{Censal ratio}

Dentro de los m�todos para proyectar poblaciones peque�as se encuentra censal-ratio (raz�n-censal). �ste m�todo se vale una variable sintom�tica asociada a la poblaci�n peque�a, a la que se supone es directamente proporcional. Un ejemplo de variable sintom�tica es la cantidad de nacimientos en una ciudad, con respecto a su poblaci�n total; otro ejemplo es la cantidad de ni�os enrolados en un curso, con respecto a un rango etario en determinado espacio geogr�fico.

El censal-ratio supone que la relaci�n entre la variable sintom�tica y la poblaci�n a estimar es una variable aleatoria cuyo valor esperado no cambia a trav�s del tiempo. 

Cuando se cuenta con estimaciones independientes de la poblaci�n total en periodos inter-censales se realiza una correcci�n sobre las proyecciones de cada �rea, tal que la suma de las proyecciones de �reas peque�as sea igual a la estimada independientemente. La correcci�n mantiene la proporci�n entre proyecciones de �reas peque�as.

Un uso frecuente de los m�todos de proyecci�n para poblaciones peque�as es realizar interpolaci�n en periodos no censales. Uno de los inconvenientes mencionados es la dificultad de realizar extrapolaci�n, puesto que para los periodos post-censales, los indicadores no est�n necesariamente disponibles. Una posible soluci�n es realizar la estimaci�n usando periodos rezagados o utilizar la �ltima proyecci�n de la variable sintom�tica como proyecci�n. \cite{Swanson2012}. 

Siguiendo a Swanson \cite{Swanson2012} tenemos: 

$$ R_{A,t,j} = \frac{S_{A,t,j}}{P_{A,t}}$$

$ R $ = Raz�n entre variable sintom�tica y poblaci�n peque�a

$ A $ = A-�sima �rea peque�a

$ S $ = Valor de la variable sintom�tica

$ t $ = periodo

$ j $ = indicador

Si tenemos un R observado en el periodo T tendremos que la proyecci�n del periodo $t + k$ ser�:

$$ \hat{P}_{A,t+k} = \frac{S_{A,t + k ,j}}{R_{A,T,j}} $$

Cuando se tiene un estimado independiente para la poblaci�n padre se realiza un control, de manera que la suma de proyecci�n para las �reas peque�as sea igual a la proyecci�n del �rea padre: 

\begin{equation}
\hat{Q}_{A,t+k} = \frac{P_{A,t+k,j}}{\sum{P_{A,t+k,j}}} * P_{t+k}
\label{equ:inpec_edad_generotasa2}
\end{equation}

$\hat{Q}_{A,t+k}$ = Proyecci�n corregida para el �rea A 

$ P_{t+k}$ = Proyecci�n independiente total en el periodo $t + k$



La variable sintom�tica cumple la tarea de agregar informaci�n conocida al modelo para hacerlo m�s preciso. En un primer paso se considera la poblaci�n privada de la libertad una poblaci�n peque�a, dentro de la poblaci�n nacional. Dentro de las estad�sticas publicadas periodicamente por el INPEC no se han encontrado variables sintom�ticas que cubran m�s de dos a�os \cite{INPEC2018}, raz�n por la cual se realiza el siguiente proceso: se utiliza una constante en el lugar de la variable sintom�tica, lo que da lugar a una proyecci�n constante por edad genero; luego ajustamos la proyecci�n al total nacional por rango de edad proyectado por el DANE. Este modelo es equivalente a proyectar una tasa de encarcelamiento constante para cada rango de edad. Para el periodo 2017 - 2020 se utiliza la tasa de encarcelamiento.

Aunque no conocemos el comportamiento de la tasa de encarcelamiento espec�fica por edad, conocemos la poblaci�n carcelaria total y la tasa de encarcelamiento. Usamos el supuesto de proporci�n estable por rango etario, para ajustar la proyecci�n 2005 - 2016 a la poblaci�n carcelaria total.

$R_e = \frac{1} {P_{c,e, 2017}} $

$\hat{P}_{c, 2017 + k} = P{c, e, 2017}$

y 

$\hat{Q}_{c, e, 2017+k} = (\hat{P}_{c, e, 2017} * P_{e,2017}) * P_{e,2017 + k}$ 

donde: 

$\hat{Q}_{c, e, 2017+k}$ = Proyecci�n de poblaci�n carcelaria ajustada, para el rango etario, en el periodo $2017 + k$ 

$P_{c, 2017}$ = Poblaci�n carcelaria en 2017, a mitad de periodo, en el rango etario e.

$P_{e,2017}$ = Poblaci�n estimada 2017 DANE a mitad de periodo en el rango etario e. 

$P_{e,2017 + k}$ = Poblaci�n estimada en $2017 + k$, del DANE a mitad de periodo en el rango etario e.

Al aplicar esta proyecci�n a periodos anteriores a 2017 tenemos una estimaci�n de lo que habr�a sido la poblaci�n carcelaria, en un escenario de tasas de encarcelamiento espec�ficas constantes. Al comparar esta proyecci�n con la poblaci�n privada de la libertad observada, es posible evidenciar el impacto que ha tenido el proceso de cambio demogr�fico en la poblaci�n carcelaria, adem�s de poner en evidencia que no es el principal factor detr�s de este cambio. El incremento debe adem�s corresponder con un incremento de las tasas de encarcelamiento espec�ficas.

Para ajustar la proyecci�n de las tasas de encarcelamiento espec�ficas se trata la poblaci�n carcelaria como un �rea padre y a las poblaciones por rango etario como poblaciones peque�as. Usando la ecuaci�n \ref{equ:inpec_edad_generotasa2} tenemos: 

$$\hat{V}_{c, e, 2017+k} = (\hat{P}_{c, e, 2017} * \hat{P}_{c,2017+k}) * P_{c,2017 + k}$$

donde: 

$\hat{V}_{c, e, 2017+k}$ = Poblaci�n carcelaria ajustada por edad. 

$ \hat{P}_{c,2017+k}$ = Poblaci�n carcelaria estimada por genero.

$ P_{c,2017 + k} $ = Poblaci�n carcelaria observada. 

\section{Poblaci�n privada de la libertad por rango de edad Febrero 2016 a Febrero 2018}

El INPEC publica un reporte mensual con estad�sticas detalladas de la poblaci�n carcelaria al cierre de mes. Los reportes est�n disponibles en la p�gina del INPEC desde e incluyen la composici�n de la poblaci�n privada de la libertad desde Enero 2013. La estructura de los reportes no es uniforme para todos los periodos. La composici�n por rangos etarios antes de febrero 2016 consolidaba en las categor�as: 18 a 29 a�os, 29 a 54 a�os, 55 a 64 a�os, y mayor de 65 a�os. A partir de marzo de 2016 se agrupa 18 a 24 y grupos quinquenales hasta los 65 a�os. \cite{INPEC2018}

En la figura \ref{fig:inpec_edad_genero2} podemos observar la composici�n demogr�fica de la poblaci�n privada de la libertad en Junio 2017. Es una poblaci�n compuesta mayoritariamente por hombres. Tanto en hombres como mujeres se observa un patr�n similar en la estructura etaria, el punto m�s alto se encuentra entre 25 y 29 a�os, 18 a 25 y 29 a 35 tienen un orden de magnitud comparable. Despu�s de los 35 a�os la poblaci�n se reduce en cada grupo quinquenal. 

En la figura \ref{fig:inpec_edad_genero3} se presenta la poblaci�n privada de la libertad, separada por genero y edad desde marzo 2016 hasta marzo 2018. El periodo de an�lisis se escogi� a conveniencia, considerando la disponibilidad de datos del INPEC con grupos etarios uniformes. La figura \ref{fig:inpec_edad_genero3} muestra que la poblaci�n carcelaria mantiene la forma de campana, con pico en 25-29, en todos los meses analizados, tanto en hombres como mujeres. En hombres el rango de edad con mayor variabilidad es el rango 18-24. En mujeres la variabilidad se presetnta en los rangos 18-24, 25-29 y 30-34. \cite{Ram2018}


\begin{figure}[H]
	\centering
	\includegraphics[width=0.7\linewidth]{gginpec_edad_genero2-1.pdf}
	\caption[Piramide poblacional, poblaci�n privada de la libertad, Junio - 2017]{Piramide poblacional, poblaci�n privada de la libertad, Junio - 2017}{Fuente: INPEC} {Elaboraci�n propia}
	\label{fig:inpec_edad_genero2}
\end{figure}


\begin{figure}	
	\centering
	\includegraphics[width=0.7\linewidth]{inpec_edad_genero3-1.pdf}
	\caption[Poblaci�n privada de la libertad. Febrero-2016 a Febrero-2018]{Poblaci�n privada de la libertad. Febrero-2016 a Febrero-2018}{Fuente: INPEC} {Elaboraci�n propia}
	\label{fig:inpec_edad_genero3}
\end{figure}

Esta variabilidad podr�a estar asociada a un componente estacional o a un componente aleatorio. Aunque la longitud de la serie disponible no permite realizar inferencia sobre componentes de estacionalidad y tendencia, la figura muestra que: A trav�s del a�o el grupo 18-24 crece, para caer abruptamente en enero del a�o siguiente; comportamiento repetido en enero 2016 y enero 2017. El comportamiento es similar en mujeres una vez se controla por el orden de magnitud. El comportamiento en las dem�s series es similar, se reduce durante todo el a�o, para subir nuevamente en enero. Este comportamiento es coherente en hombres y mujeres, una vez se controla la magnitud. \ref{fig:inpec_edad_genero3} Este comportamiento implica que no es conveniente trabajar la proyecci�n de la tasa de encarcelamiento a nivel mensual pues no se puede estimar el aparente componente de estacionalidad de la serie. Por conveniencia se trabaja con la proyecci�n de poblaci�n a mitad de periodo, y se trata la tasa de encarcelamiento como una variable aleatoria cuyo valor esperado se mantiene a trav�s del tiempo.

\begin{figure}[H]
	\centering
	\includegraphics[width=0.7\linewidth]{inpec_edad_genero1-1.pdf}
	\caption[Poblaci�n privada de la libertad. Febrero-2016 a Febrero-2018]{Poblaci�n privada de la libertad por rango etario. Febrero-2016 a Febrero-2018}{Fuente: INPEC} {Elaboraci�n propia}
	\label{fig:inpec_edad_genero1}
\end{figure}

Puesto que no se cuenta con informaci�n de periodos m�s extensos se valida an�lisis a un menor nivel de granularidad. Dentro de la informaci�n disponible se encuentra separada la poblaci�n carcelaria por regional. En la figura \ref{fig:inpec_edad_generoregion2} se observa nuevamente, en todas las regionales, la poblaci�n carcelaria con forma de campana. Consistentemente las mujeres tienen mayor variabilidad que los hombres en cada rango de edad y poblaciones carcelarias m�s peque�as se encuentran asociadas a una mayor variabilidad.

La forma del aparente componente estacional tambi�n se replica a nivel regional. Ver figura \ref{fig:inpec_edad_generoregion3}. 

\begin{figure}[H]
	\centering
	\includegraphics[width=1.1\linewidth]{inpec_edad_generoregion2-1.pdf}
	\caption[Poblaci�n privada de la libertad. Febrero-2016 a Febrero-2018]{Poblaci�n privada de la libertad por rango etario, regi�n. Febrero-2016 a Febrero-2018}{Fuente: INPEC} {Elaboraci�n propia}
	\label{fig:inpec_edad_generoregion2}
\end{figure}

\begin{figure}[H]
	\centering
	\includegraphics[width=1.1\linewidth]{gginpec_edad_genero_region-1.pdf}
	\caption[Poblaci�n privada de la libertad. Febrero-2016 a Febrero-2018]{Poblaci�n privada de la libertad por rango etario, regi�n. Febrero-2016 a Febrero-2018}{Fuente: INPEC} {Elaboraci�n propia}
	\label{fig:inpec_edad_generoregion3}
\end{figure}

Las personas no aparecen de la nada en la carcel. 


\begin{figure}[H]
	\centering
	\includegraphics[width=0.7\linewidth]{gginpec_edad_generotasa1-1.pdf}
	\caption[Tasa de espec�fica de encarcelamiento. Junio-2017]{Tasa de espec�fica de encarcelamiento. Junio-2017}{Fuente: INPEC, DANE} {Elaboraci�n propia}
	\label{fig:inpec_edad_generotasa1}
\end{figure}

\begin{figure}[H]
	\centering
	\includegraphics[width=0.7\linewidth]{inpec_edad_generotasa2-1.pdf}
	\caption[Tasa espec�fica de encarcelamiento. Junio-2016, Junio -2017]{Tasa espec�fica de encarcelamiento. Junio-2016, Junio -2017}{Fuente: INPEC, DANE} {Elaboraci�n propia}
	\label{fig:inpec_edad_generotasa2}
\end{figure}

\section{Proyecciones Censal Ratio}

\begin{figure}[H]
	\centering
	\includegraphics[width=0.9\linewidth]{Proyecciones_DANE-1.pdf}
	\caption[Poblaci�n privada de la libertad. 2005-2020]{Proyecci�n de la poblaci�n privada de la libertad,  2005 - 2020, ajustado por rango etario DANE}{Fuente: INPEC, DANE} {Elaboraci�n propia}
	\label{fig:inpec_edad_pry}
\end{figure}



   \chapter{Modelos estado-espacio}

\section{Selecci�n de Software}

Para mantener un �nico ambiente de desarrollo, nos enfocamos en los paquetes desarrollados en R, como los mencionados por Petris \cite{Petris2011}, Helske \cite{Helske2017} y Tusell \cite{Tusell2011}. Particularmente usaremos el paquete \textit{Kfas} por Helske(2017) \cite{Helske2017}


\section{Identificaci�n del modelo}

\section{Simulaci�n Monte-Carlo}

\section{Estimaci�n de par�metros}

\section{Proyecciones 2017 - 2020}

\section{Conclusiones}
%\appendix
%    \chapter{Gr�ficas adicionales}

\begin{figure}[htb]
 \centering
 \includegraphics[width=10cm]{small_escudo_color}
 \caption{Escudo oficial de la UN a color dise�ado por el Maestro Francisco Duarte.}
 \label{fig:escudo_color}
\end{figure}


\begin{table}
\centering
\caption{Unidades de \TeX.}
\begin{tabular}{l|l}\hline
\verb"mm" & mil�metro $\approx$ 1/25 pulgada \\
\verb"cm" & cent�metro = 10 mm \\
\verb"in" & pulgada $\approx$ 25 mm \\
\verb"pt" & punto $\approx$ 1/72 pulgada $\approx$ 1/3 mm \\
\verb"em" & aprox. el ancho de una m en el tipo actual \\
\verb"ex" & aprox. la altura de una x en el tipo actual \\\hline
\end{tabular}
\end{table}

\begin{table}
\centering\renewcommand{\arraystretch}{2}
\caption{Tama�os de los tipos de fuentes \LaTeX.}
 \begin{tabular}{ll|ll}\hline
{\tiny\verb"\tiny"} & {\tiny letra diminuta}
& {\large\verb"\large"} & {\large letra grande} \\
{\scriptsize\verb"\scriptsize"} & {\scriptsize letra muy peque�a}
& {\Large\verb"\Large"} & {\Large letra mayor} \\
{\footnotesize\verb"\footnotesize"} & {\footnotesize letra bastante peque�a}
& {\LARGE\verb"\LARGE"} & {\LARGE muy grande} \\
{\small\verb"\small"} & {\small letra peque�a}
& {\huge\verb"\huge"} & {\huge enorme} \\
\verb"\normalsize" & letra normal
& {\Huge\verb"\Huge"} & {\Huge la mayor} \\\hline
 \end{tabular}
\end{table}





\backmatter
    \conclusion{conclusiones}
    \futurework{trabajofuturo}
%    \glossary{glosario}
%\nocite{*}
    \bibliography{references}
    %\printindex
\end{document}
